\section{Bayesian Contract}

The implemented record is an \path{AuditAlgorithmRun} whose method id is
\texttt{bayesian-tabulation-audit-v1}. Its Bayesian fields are explicit and are
kept out of non-Bayesian audit methods.

\begin{center}
\small
\begin{tabularx}{\textwidth}{lX}
\toprule
Field & Meaning \\
\midrule
\texttt{method\_id} & Versioned method name. Current value: \texttt{bayesian-tabulation-audit-v1}. \\
\texttt{bayesian\_prior\_id} & Non-empty identifier for the declared prior specification. \\
\texttt{bayesian\_likelihood\_id} & Non-empty identifier for the declared likelihood or evidence model. \\
\texttt{posterior\_winner\_probability\_ppm} & Declared posterior probability, in parts per million, that the reported winner remains the winner. \\
\texttt{posterior\_risk\_ppm} & Declared posterior upset probability, in parts per million. \\
\texttt{simulation\_seed} & Optional seed metadata for a posterior computation performed outside this verifier. \\
\texttt{posterior\_draws} & Optional positive count of posterior draws. \\
\texttt{calibrated\_risk\_limit\_ppm} & Optional declared calibration threshold; the current verifier only checks its shape. \\
\texttt{assertions} & Non-empty list whose assertion kinds must be \texttt{bayesian-outcome}. \\
\bottomrule
\end{tabularx}
\end{center}

The contract is intentionally a record-shape and design boundary. The prior and
likelihood identifiers make the posterior claim attributable. The posterior
winner and upset-risk fields are bounded probability fields. The optional
calibration field records a possible bridge to risk-limiting language, but it
is not treated as a valid frequentist risk limit unless a later calibrated rule
is implemented.
